\section{Declarações de integridade científica e ciência aberta}

\subsection{Conflitos de interesse e financiamento}

Os autores declaram não possuir conflitos de interesse relacionados ao desenvolvimento, à avaliação ou à apresentação dos resultados deste trabalho. A prova de conceito foi desenvolvida em contexto acadêmico, sem financiamento externo específico declarado para esta pesquisa.

\subsection{Disponibilidade de dados, código e artefatos}

Os dados utilizados derivam de registros públicos de discursos parlamentares do Senado Federal e foram organizados em \textit{dataset} público no Hugging Face\footnote{\href{https://huggingface.co/datasets/fabriciosantana/discursos-senado-legislatura-56}{\textit{Dataset} público de discursos da 56\textordfeminine{} Legislatura}.}. O código, os \textit{scripts} auxiliares, os parâmetros de preparação da base, a rubrica, o \textit{prompt} operacional e os artefatos experimentais citados no artigo estão disponíveis no repositório público do projeto\footnote{\href{https://github.com/fabriciosantana/mcdia/tree/main/05-iag/4-project}{Repositório público do projeto}.}. Essa disponibilização busca favorecer auditoria, reprodutibilidade e reutilização crítica do protocolo em estudos posteriores.

\subsection{Uso de inteligência artificial generativa}

Ferramentas de inteligência artificial generativa foram utilizadas como apoio à programação, à revisão textual, à organização de argumentos e à avaliação automatizada descrita na metodologia. As decisões de escopo, interpretação dos resultados, seleção das evidências, redação final e responsabilidade pelo conteúdo permanecem sob responsabilidade dos autores. As saídas geradas por modelos foram revisadas criticamente e não substituem validação humana independente.

\subsection{Considerações éticas}

O estudo utiliza documentos parlamentares públicos, sem coleta direta de dados junto a participantes humanos e sem acesso a informações sigilosas. Ainda assim, por tratar de acervo institucional e de respostas geradas por IA, o trabalho adota cautela na interpretação dos resultados e enfatiza a necessidade de verificação documental, supervisão humana e governança adequada antes de qualquer uso operacional em contexto público.
